\documentclass[11pt]{article}
% ======================================================================
%  Preambule commun - Sujets Bac Serie A (Mathematiques)
%  Style sobre : noir et blanc, sans couleurs, sans filigrane,
%  sans en-tetes ni pieds de page decoratifs.
% ======================================================================
\usepackage[utf8]{inputenc}
\usepackage[T1]{fontenc}
\usepackage{lmodern}
\usepackage{textcomp}
\usepackage{amsmath,amssymb}
\usepackage{enumitem}
\usepackage{array}
\usepackage[a4paper,margin=2.3cm]{geometry}
\usepackage{fancyhdr}
\usepackage{tikz}

% Symbole degre robuste + justification souple
\DeclareUnicodeCharacter{00B0}{\ensuremath{^\circ}}
\tolerance=2000
\emergencystretch=2em
\frenchspacing

% En-tete / pied de page minimalistes (numero de page seul)
\pagestyle{fancy}
\fancyhf{}
\renewcommand{\headrulewidth}{0pt}
\fancyfoot[C]{\thepage}

% Macros mathematiques
\newcommand{\R}{\mathbb{R}}
\newcommand{\N}{\mathbb{N}}
\newcommand{\Z}{\mathbb{Z}}
\newcommand{\vi}{\vec{\imath}}
\newcommand{\vj}{\vec{\jmath}}

% Titres de blocs
\newcommand{\exo}[2]{\bigskip\noindent{\large\bfseries Exercice #1}\hfill\textit{#2}\par\nopagebreak\smallskip}
\newcommand{\partie}[1]{\medskip\noindent{\bfseries Partie #1}\par\nopagebreak\smallskip}
\newcommand{\entetesujet}[1]{\begin{center}{\Large\bfseries #1}\end{center}\vspace{0.3em}\hrule\vspace{1.1em}}

% Questions numerotees : niveau 1 = chiffres gras, niveau 2 = a. b. c.
\newlist{questions}{enumerate}{1}
\setlist[questions]{label=\textbf{\arabic*.},leftmargin=1.8em,itemsep=3pt,topsep=3pt}
\newlist{sousq}{enumerate}{1}
\setlist[sousq]{label=\textbf{\alph*.},leftmargin=1.8em,itemsep=2pt,topsep=2pt}


\begin{document}

\entetesujet{Sujet bac 2020 -- Série A}

\exo{1}{8 points}
\partie{1}
\begin{questions}
  \item Répondre par vrai ou faux à chacune des affirmations suivantes :
  \begin{sousq}
    \item Chacun des nombres 28 et 56 est un multiple de 7.
    \item L'équation $6x^2 - 7x + 3 = 0$ a 2 solutions réelles distinctes $x'$ et $x''$.
    \item Le multiple d'un nombre $a$ est le produit de $a$ par un entier $n$.
    \item Le nombre 124 est divisible par 3.
  \end{sousq}
  \item \begin{sousq}
    \item Rendre irréductible la fraction $\dfrac{1950}{235}$.
    \item Trouver un encadrement du nombre $\dfrac{1950}{235}$ à $10^{-3}$ près.
  \end{sousq}
\end{questions}

\partie{2}
\begin{questions}
  \item Résoudre dans $\R$ l'équation : $\dfrac{2x - 1}{x + 4} = \dfrac{2x + 1}{x - 2}$.

  \underline{NB.} Préciser les conditions d'existence. \underline{Indication :} le produit des moyens est égal au produit des extrêmes.
  \item \begin{sousq}
    \item Résoudre dans $\R$ l'équation $(E) : 3x^2 - 2x - 1 = 0$.
    \item En déduire les solutions de l'équation $(E') : 3\,e^{2x} - 2\,e^{x} - 1 = 0$.
  \end{sousq}
  \item \begin{sousq}
    \item Résoudre dans $\R^2$ le système $S : \begin{cases} x - 2y = -1 \\ 2x + 3y = -2 \end{cases}$
    \item En déduire les solutions du système $(S')$ suivant : $\begin{cases} \ln x - 2\ln y = -1 \\ 2\ln x + 3\ln y = -2 \end{cases}$
  \end{sousq}
\end{questions}

\exo{2}{8 points}
On considère la fonction numérique $g$ de la variable réelle $x$ définie telle que $g(x) = \dfrac{-2x + 2}{2x + 3}$. On désigne par $(\mathcal{C})$ la courbe représentative de la fonction $g$ dans un plan muni d'un repère orthonormé $(O,\vi,\vj)$ d'unité graphique 1 cm.
\begin{questions}
  \item \begin{sousq}
    \item Calculer l'image du réel $-1$.
    \item Justifier que le réel $-\dfrac{3}{2}$ n'a pas d'image.
  \end{sousq}
  \item Calculer la limite de $g$ en $+\infty$ et en $-\dfrac{3}{2}$ à gauche et à droite.
  \item \begin{sousq}
    \item Déterminer la dérivée $g'$ de $g$ pour tout réel $x \ne -\dfrac{3}{2}$.
    \item Préciser le signe de $g'(x)$ pour $x \in \left]-\infty\,;-\dfrac{3}{2}\right[ \cup \left]-\dfrac{3}{2}\,;+\infty\right[$.
  \end{sousq}
  \item Dresser le tableau de variation de $g$ sur $\left]-\infty\,;-\dfrac{3}{2}\right[ \cup \left]-\dfrac{3}{2}\,;+\infty\right[$. On donne la valeur $-1$ à la limite de $g$ en $-\infty$ ($\lim\limits_{x\to-\infty} g(x) = -1$).
  \item Recopier et compléter le tableau de valeurs suivant :
  \begin{center}\renewcommand{\arraystretch}{1.5}
  \begin{tabular}{|c|c|c|c|c|}
  \hline
  $x$ & $-2$ & & $0$ & $1$ \\
  \hline
  $g(x)$ & & $4$ & & \\
  \hline
  \end{tabular}
  \end{center}
  \item Tracer $(\mathcal{C})$ et les droites d'équations $x = -\dfrac{3}{2}$ et $y = -1$ dans le plan muni du repère $(O,\vi,\vj)$.
\end{questions}

\exo{3}{4 points}
Dix adolescents s'exercent à lancer le poids du bras droit puis du bras gauche. Les résultats (distances en mètres) obtenus sont consignés dans le tableau suivant :
\begin{center}\renewcommand{\arraystretch}{1.4}
\begin{tabular}{|c|c|c|c|c|c|c|c|c|c|c|}
\hline
Bras droit ($x_i$) & 5{,}5 & 7{,}1 & 5{,}8 & 6{,}4 & 6 & 6{,}2 & 7{,}2 & 5{,}6 & 6{,}8 & 5{,}6 \\
\hline
Bras gauche ($y_i$) & 4{,}1 & 6{,}2 & 4 & 5{,}5 & 4{,}9 & 4{,}7 & 6 & 4{,}9 & 5 & 3{,}5 \\
\hline
\end{tabular}
\end{center}
\begin{questions}
  \item \begin{sousq}
    \item Reproduire et compléter le tableau $(T)$ suivant :
    \begin{center}\renewcommand{\arraystretch}{1.5}
    \begin{tabular}{|c|c|c|c|c|c|c|c|c|c|c|}
    \hline
    $x_i$ & 5{,}5 & 5{,}6 & 5{,}6 & 5{,}8 & 6 & 6{,}2 & 6{,}4 & 6{,}8 & 7{,}1 & 7{,}2 \\
    \hline
    $y_i$ & & & & & & & & & & \\
    \hline
    \end{tabular}
    \end{center}
    \item On considère le nuage formé par les cinq (5) premiers points du tableau $(T)$ ci-dessus. Calculer les coordonnées du point moyen $G_1(\overline{x}_1, \overline{y}_1)$ de ce nuage (arrondir à 2 chiffres après la virgule).
    \item On considère le nuage formé par les cinq (5) derniers points du tableau $(T)$ ci-dessus. Calculer les coordonnées du point moyen $G_2(\overline{x}_2, \overline{y}_2)$ de ce nuage (arrondir à 2 chiffres après la virgule).
  \end{sousq}
  \item \begin{sousq}
    \item Montrer que l'équation de la droite passant par les points moyens $G_1$ et $G_2$ est : $y = 1{,}077x - 1{,}778$.
    \item Pour un adolescent qui lance le poids à 4,4 m avec le bras gauche, estimer la distance du lancer de poids de son bras droit.
  \end{sousq}
\end{questions}

\end{document}
